\section{\MakeUppercase{Introduction}}
The aircraft maintenance is extremely important and can be defined as "the process of ensuring that a system continually performs its intended function at its designed-in level of reliability and safety." \textcite[p. 35]{kinnison2013aviation}. It is possible to optimize the aircraft availability by proceeding in a preventive way, avoiding delays, excessive cost ans even catastrophic damages. Those proceeding can be structured based on RCM or "\textit{Reliability-Centered Maintenance}", first developed in the 70's for the U.S military, this process was based in civilian aviation and It's first priority is the preventive analysis of system's integrity instead of reactive procedures to damage \cite[p. 34]{kinnison2013aviation}. Among the schedule on-condition tasks  described by the RCM, crack length monitoring is used to define a potential failure point and crack growth rate to establish an inspection interval \cite[p. 52]{nowlan1978maintenance}. In modern aircraft industry, the RCM can be used both in development and operation phases \cite{ludek2005rcm}, and is right at the development phase that is possible to improve reliability, because at the operation all that is left to be done is to repair or replace. Therefore, It is necessary to identify the failure modes before defining the points to be improved. The failure can be defined as when the component can no longer accomplish the goal for what It was designed \cite{findlay2002aircraft} and while the most common failure mode of general engineering components is the corrosion, aircraft parts are mainly affected by fatigue \cite{tavares2017overview}. In order to improve the fatigue life of a component, changing the geometry to avoid stress concentration points may be effective, but not always possible due project restrictions, for example: a fastener hole causes stress concentration but It's shape can not be changed and neither does the fastener position. Thus, the project must have certain level of damage tolerance so that new operation conditions or replacement can be established. As the crack growth rate can define the inspection period, reducing this rate can increase the fatigue life and inspection window of the component. So, methods such as surface treatments that introduces compressive stress along the thickness of the material can be good alternatives for reducing the tractive stress at the crack tip. One of these method is the  Laser Shock Peening [...]. Experiments using $5 mm$ thickness sheet of Ti treated with LSP done by \textcite{SUN2021106081} showed crack closing effect and a reduction in crack growth rate, this study proposes a mathematical model that can reproduce the crack growth curves under LSP in $2 mm$ sheet of Al-5087.
